\begin{resumo}[Abstract]
 \begin{otherlanguage*}{english}
  This work aims to propose a comparative analysis between two processing methods for \gls{smart_metering} systems:
  \gls{edge_computing}, with \gls{ocr} inference performed locally on the \gls{esp32cam} microcontroller through a 
  supervised trained machine learning model, and \gls{cloud_computing}, with the transmission of the captured image 
  for remote processing by the Tesseract engine combined with the OpenCV library. The motivation lies in the high cost of replacing
  various conventional meters, widely used in residential and private environments, with native smart models, which 
  favors \textit{\gls{retrofitting}} as a more economically viable alternative. Data transmission between edge nodes and the central
  gateway is carried out via a \gls{lorawan} network using Spreading Factor SF8, a configuration that ensures compliance with the 
  regulatory dwell time limit of 400~ms established by ANATEL's Act No.~14.448, while maximizing reach within this limit. 
  The feasibility of sending fragmented images via \gls{lorawan} was analytically demonstrated, where a 160$\times$120 pixel 
  grayscale image in JPEG format of approximately 3.10~KB can be transmitted in 38 packets in about 13.3~s of active radio time. 
  The two processing workflows will be evaluated according to three quantitative metrics: \gls{ocr} reading accuracy (\%), end-to-end latency (ms),
  and energy consumption per reading (mAh). The results are expected to guide the choice of the most suitable architecture 
  for large-scale \textit{Smart Metering} deployments, considering the trade-off between energy efficiency, recognition accuracy,
  and system maintenance complexity.
  \vspace{\onelineskip}
  \noindent
  \textbf{Keywords}: Smart Metering. ESP32. LoRaWan. OCR.
 \end{otherlanguage*}
\end{resumo}